\section{Executable reference corpora and comparison design}
\label{sec:evaluation}

The executable references provide reproducible observations of communication
and exposure. The pinned SPQR source is
\texttt{f2589fef855c10f39d72634dab3d14654dd410bf}. Seven schedules each contain
2,048 sends: alternating, loss, reordering, duplication, 9:1 asymmetry,
offline-then-exchange, and one-way traffic. Six capture points for each endpoint
produce 84 passive snapshot cases. The whole-KEM control runs those schedules
at local-send intervals 1, 32 and 64, producing 252 corresponding cases.
Together these are 28 configuration/schedule combinations and 336 snapshots.
The explicitly modified 64-byte SPQR variant has its own source identity.

\subsection{Observed derivation and communication}

The snapshot predictor receives one endpoint's serialized state and the full
public transcript, including packets lost before peer delivery. A separate
harness compares every derived key with the actual sender's key. Public
deterministic seeds are test fixtures; the measurement establishes causal key
derivation through the declared predictor interface. The output label
\texttt{not\_derived} denotes that predictor's result. It is separate from the
security predicate of a computational experiment.

Table~\ref{tab:reference} reproduces the retained alternating-schedule
observations. The snapshot is taken at $A$ immediately after its first send,
leaving 2,047 future sends. The pending KEM key accounts for 87 of the 172
derived keys in the pinned 32-byte SPQR case, demonstrating why the currently
emitted chain epoch alone does not describe exposure.

\begin{table}[t]
\centering
\small
\begin{tabular}{@{}lrrrr@{}}
\toprule
Component/profile & Mean bytes/send & Peak frame & Epochs $A/B$ & Derived keys\\
\midrule
Pinned SPQR, 32 B & 32.150 & 37 & 24/23 & 172\\
SPQR variant, 64 B & 58.971 & 69 & 46/45 & 88\\
Whole KEM, interval 1 & 819.658 & 1,225 & 683/683 & 4\\
Whole KEM, interval 32 & 44.047 & 1,223 & 32/32 & 66\\
Whole KEM, interval 64 & 25.023 & 1,223 & 16/16 & 130\\
\bottomrule
\end{tabular}
\caption{Retained component-level alternating workload. Byte counts are the
component wire contribution; epoch counts are each component's root/output
observations. The constructions have different authentication and retention
graphs. Outer application AEAD, transport, durable storage and energy are
outside these measurements.}
\label{tab:reference}
\end{table}

The one-way schedule produces no fresh PQ epoch in these references and the
predictor derives every post-snapshot key. The whole-KEM active-control
experiment addresses a different property: one initial endpoint-state capture
allows two attacker forks to replace subsequent exchanges while honest
confirmed epochs increase. All 1,536 sender keys across three 512-send profiles
match the attacker-derived keys. This supplies a concrete distinction between
fresh KEM execution and authenticated recovery when the ongoing authentication
authority is itself derived from captured session state.

\subsection{Aligning the Continuity experiment}

A matched experiment must give each construction the same application and
control-message opportunities and the same loss, reorder and duplication
schedule. Authentication capabilities are recorded explicitly: the whole-KEM
control uses session-state authentication, whereas Continuity controls use
the separately modeled device signing authority. Serialized volatile state,
retained reservations and actual durable journal images are distinct disclosure
objects.

The planned Continuity comparison records, for each exposure case, the original
generation and encapsulation reservation identifiers, the retained record
projection, acquired wrapping-key identifiers, endpoint phase/counter state,
public-wire availability and the actually reconstructed message set.
Communication records include total bytes and largest frame. Persistence
records include durable writes, restart/retry work and retained-state size.
Each timing denominator covers its declared operation interval and names the
compiler profile, primitive backend, host and sample unit.

The principal ablations hold the underlying reservation fixed while changing
its execution time, delay a receipt, copy a target through a sibling aggregate
write, and compare origin-preserving labels with labels based on the last image
write. These ablations test the mechanisms isolated by the theorem and
separating executions. Their comparison will explain both exposure and cost
under each matched workload.
